% !TeX root = ../kazymyr-thesis.tex

\chapter{Preliminaries}\label{chapter:preliminaries}

\section{Transactions and concurrency control}\label{sec:prelim-transactions}

\subsection{Transactions and serialisability}\label{sec:prelim-serialisability}

A transaction is a partial order of read and write steps over resources. A
schedule of concurrently executing transactions is correct if it is
equivalent to some serial execution of the same
transactions~\cite{weikum2001transactional}. This is the
only correctness criterion used in the thesis; the vocabulary is needed
mainly to state what the ACID execution path guarantees and what the BASE
path gives up.

\subsection{ACID and two-phase locking}\label{sec:prelim-acid}

The simulator enforces isolation with conservative two-phase locking and
cautious waiting~\cite{hsu1992cautious}. A transaction declares the resources it needs before it
begins, acquires them as a set rather than incrementally, and releases
them only after its final step. A transaction that would have to wait for
a resource held by another does not acquire a partial lock set and hold it
while waiting.

Three consequences of this scheme are used later. Locks are never acquired
after the first release, so no transaction can extend its lock set while
holding one. A waiting transaction holds no locks that another transaction
is waiting for. Therefore the wait-for relation among ACID transactions
cannot contain a cycle, and pure ACID execution is deadlock-free by
construction. Section~\ref{sec:approach-deadlock} shows that the arbiter
introduced in this thesis preserves that property.

\subsection{BASE and eventual consistency}\label{sec:prelim-base}

BASE trades immediate consistency for availability and
throughput~\cite{pritchett2008base,vogels2009eventually}: a transaction
commits before all of its consequences have been applied, and the system
converges afterwards. Two terms recur throughout the thesis. A
resource is \emph{dirty} when a consistency-preserving update that
concerns it has not yet been applied. \emph{Reconciliation} is the
background work that applies those pending updates and clears the
resources again.

\section{Multi-model systems and consistency preservation}\label{sec:prelim-multimodel}

\subsection{Models, semantic overlap and CPRs}\label{sec:prelim-cprs}

A model is a set of resources describing one view of the
product~\cite{klare2021vitruvius}. Two
models exhibit \emph{semantic overlap} when they contain resources
representing the same underlying fact; the overlap degree is a
configuration parameter of the simulator. A consistency preservation rule
(CPR) states that a write to one resource requires updates to the
resources that overlap it. Rules compose: an update produced by a rule may
itself trigger further rules, and the resulting chain is bounded by a
maximum CPR depth.

The transactions that carry this induced work are called \emph{bookkeeper}
transactions. They are spawned by the root transaction whose write
triggered the rule and inherit its execution mode. A rule also has a
referential action: \texttt{ON\_CASCADE} propagates the update, whereas
\texttt{ON\_RESTRICT} rejects the originating write instead. All
experiments in this thesis use \texttt{ON\_CASCADE}.

\subsection{The dirty set as a control signal}\label{sec:prelim-dirtyset}

Under BASE semantics, the resources awaiting a CPR update form the dirty
set $D_k$, of size $|D_k|$, and the dirty fraction is $d_k = |D_k| / N$
over the $N$ resources of the system. Reconciler processors drain this set
in the background; their number $R$ is the main provisioning parameter of
the system. The dirty fraction is the quantity every result in this thesis
is about: it is what the adaptive policy observes, what the critical
threshold bounds, and what the reported percentiles summarise.

It is also the only signal available for deciding \emph{when} the system
should stop accumulating inconsistency, which makes it a control signal
and not merely a metric. The simplest policy over such a signal is a
single threshold: run BASE below it, ACID above it. A single threshold on
a fluctuating signal causes flapping, that is, the system switches back
and forth every time the signal crosses it, paying the cost of a mode
change without staying in the new mode long enough to benefit from it. The
standard remedy is a \emph{hysteresis band}: separate thresholds for
entering and for leaving a state. Here the critical threshold $c$ governs
entry into ACID and the lower recovery threshold $r$ governs the return to
BASE, so the interval $[r, c]$ is the margin the signal must traverse
before a second switch can occur. Section~\ref{sec:approach-adaptive-oracle}
states the resulting rule for the oracle built in this thesis.

Two properties of the dirty fraction make this band cheap to reason about.
It carries no measurement noise --- $d_k$ is an exact count of a set the
simulator maintains, not an estimate --- and it is bounded in $[0, 1]$, so
both thresholds are directly interpretable as fractions of the resource
population.

\section{The base simulator}\label{sec:prelim-simulator}

\subsection{Execution engine and scheduling}\label{sec:prelim-engine}

The work extends the discrete-event simulator of a transactional server
for multi-model systems introduced by Zacouris and
Acosta~\cite{zacouris2025m2ts,zacouris2026mxs}. Time advances in fixed steps of one
simulated second. Transactions are admitted, expanded into their CPR
closure, scheduled onto IO and CPU processor pools, and either commit or
abort; under BASE the closure is instead queued for the reconciler pool.
Because the clock is virtual, results do not depend on the speed of the
host machine.

\subsection{Parameters and metrics}\label{sec:prelim-parameters}

The simulator exposes the workload profile, the number and size of models,
the semantic overlap degree, the processor and reconciler counts, the lock
budget, and the maximum number of resources and models per transaction.
The workload profile decides how root transactions arrive; the closed-loop
profile \texttt{EQUILIBRIUM} used throughout this thesis holds the number
of in-flight transactions constant and admits a new one whenever one
completes, so the load is stationary by construction
(Section~\ref{sec:exp-workload}). The simulator also generates its own
workload from seeds, so no external dataset is involved, and it exports
per-run metrics as CSV. The parameter table
of~\cite{zacouris2026mxs} is reproduced here, and the configuration used
in Chapter~\ref{chapter:experiments} is stated as a difference against
it.

\section{Notation}\label{sec:prelim-notation}

Table~\ref{tab:notation} fixes the symbols used throughout. Lowercase
symbols are fractions of the resource population and uppercase symbols are
absolute counts or configuration values. The index $k$ counts successive
samples, which are taken once per admission event rather than once per
unit of simulated time --- a property with consequences of its own
(Section~\ref{sec:recovery-mechanism}).

\begin{table}[htbp]
  \centering
  \caption{Notation used throughout this thesis.}
  \label{tab:notation}
  \small
  \begingroup
  \renewcommand{\arraystretch}{1.3}
  \begin{tabular}{llll}
    \toprule
    Symbol                 & Meaning                  & Definition                      & Implementation                        \\
    \midrule
    $R$                    & Reconciler processors    & ---                             & \texttt{numberOfReconcilerProcessors} \\
    $N$                    & Total resources          & ---                             & \texttt{totalResources}               \\
    $L$                    & Max resources per txn    & ---                             & \texttt{maxTransactionResources}      \\
    $D_k$                  & Dirty resources (count)  & ---                             & \texttt{DirtySetManager.size()}       \\
    $d_k$                  & Dirty fraction           & $D_k / N$                       & \texttt{dirtyFraction}                \\
    $c$                    & Critical threshold       & ---                             & \texttt{criticalDirtyFraction}        \\
    $r$                    & Recovery threshold       & --- ($0 \le r \le c$)           & \texttt{recoveryDirtyFraction}        \\
    $p$                    & Random-oracle ACID prob. & ---                             & \texttt{randomOracleAcidProbability}  \\
    $\phi_{\mathrm{time}}$ & ACID time-share          & (Section~\ref{sec:exp-metrics}) & \texttt{acid\_fraction\_time}         \\
    $\phi_{\mathrm{comp}}$ & ACID completion-share    & (Section~\ref{sec:exp-metrics}) & \texttt{acid\_fraction\_completed}    \\
    \bottomrule
  \end{tabular}
  \endgroup
\end{table}
